\documentclass[10pt,twoside]{article}
\usepackage{amssymb,amsmath,amsthm,longtable}
\usepackage{latexsym}
\usepackage{graphicx}
\usepackage{hyperref}
\usepackage[top=1in, bottom=1in, left=1in, right=1in]{geometry}

% some definitions

\def\cross{\times}
\def\del{\nabla}
\def\grad{\vec\nabla}
\def\div{\grad\cdot}
\def\curl{\grad\cross}
\def\I{I\!\!\!-}
\def\dd{{\rm d}}
 
% begin equation, itemize, etc.

\def\be{\begin{equation}}
\def\ee{\end{equation}}
\def\bi{\begin{itemize}}
\def\ei{\end{itemize}}
\def\ben{\begin{enumerate}}
\def\een{\end{enumerate}}
\def\i{\item{}}

\def\defn{\underline{Definition}:\ }
\def\ex{\underline{Example}:\ }
\def\exer{\underline{Exercise}:\ }
\def\soln{\underline{Solution}:\ }
\def\theor{\underline{Theorem}:\ }
\def\pf{\underline{Proof}:\ }
\def\ques{\underline{Question}:\ }
\def\ans{\underline{Answer}:\ }
 
%%%%%%%%%%%%%%%%%%%%%%%%%%%%%%%%%%%%%%%%%%%%%%%%%%%%%%%%%%%%%%
            
\begin{document}

\title{College Physics II \\
Lecture summaries}
%\author{J.D.\ Romano}
\author{DRAFT}
\date{Fall 2026}
\maketitle
 
\begin{abstract}
\input{abstract}
\end{abstract}

\newpage
\tableofcontents
\cleardoublepage

\section{Electric charges, forces, and fields}
\label{s:charges_etc}

\subsection{Electric charge} 

\bi

\i Properties of electric charge:

(i) Two types: {\em positive} and {\em negative}
($+$ and $-$), whose effects tend to cancel.
Opposite charges attract, like charges repel.

(ii) Total electric charge is {\em conserved}
(cannot be created or destroyed, only transferred from one 
object to another).

(iii) Electric charge is {\em quantized}
(integer multiples of the
charge of a proton $e = 1.602\times 10^{-19}~{\rm C}$;
an electron has charge $-e$).

\i Materials:

(i) {\em Conductors}: one or more electrons per atom are able to 
move freely through the material.
Examples: aluminum, copper, gold, ... (metals).

(ii) {\em Insulators}: electrons are bound to atomic nuclei.
Examples: rubber, wood, paper, ...

\i Some demos: 

1) A plastic rod rubbed with fur is negatively charged;
a glass rod rubbed with silk is positively charged.
Show that opposite charges attract, like charges repel.

%2) An uncharged aluminum can will be attracted to both
%negative and positive rods.
%Electrons in the aluminum can (a conductor) move away 
%from the negative rod or toward the positive rod 
%leading to an attractive force in both cases.

2) Bits of uncharged paper will be attracted to both
negative and positive rods due to charge separation 
within the atoms / molecules in the paper (an insulator).

3) A charged rubber ballon sticks to the white board
for the same reason as (2).

\ei

%%%%%%%%%%%%%%%%%%%%%%%%%%%%%%%%%%%%%%%%%%%%%%%%%%%%
\subsection{Electric force}
\bi

\i {\em Coulomb's law}:
The magnitude of the electric force between two point 
charges $q_1$ and $q_2$ separated by a distance $r$ is
%
\be
F = \frac{k |q_1| |q_2|}{r^2}
\quad{\rm where}\quad
k= 8.99\times 10^9~
\frac{{\rm N}\cdot{\rm m}^2}{{\rm C}^2}
\ee
%

\i The force is attractive (directed towards one another)
if the two charges have opposite signs;
the force is repulsive (directed away from one another)
if the two charges have the same signs.

\i One sometimes writes
%
\be
k=\frac{1}{4\pi\epsilon_0}\,,
\quad{\rm where}\quad
\epsilon_0= 8.85\times 10^{-12}~\frac{{\rm C}^2}{{\rm N}\cdot{\rm m}^2}
\ee
where $\epsilon_0$ 
is the so-called {\em permitivity} of free space (vacuum).

\i The electric force on point charge $Q$ due to a set of 
point charges $q_1$, $q_2$, ..., $q_N$ is given by the vector 
sum of the forces 
%
\be
\vec F_{Q,{\rm net}} = 
\vec F_{1Q} + \vec F_{2Q} + \cdots + \vec F_{NQ}
\ee
%
where $\vec F_{1Q}$ is the electric force that point charge $q_1$ 
exerts on $Q$, etc.
This is called the {\em superposition principle}.

\ei

%%%%%%%%%%%%%%%%%%%%%%%%%%%%%%%%%%%%%%%%%%%%%%%%%%%%
\subsection{Electric field}

\bi

\i The electric force between charges is a {\em non-contact}
force.
A set of charges creates a {\em field} to which another 
charge responds.

\i The electric field at a point $P$ in space is a vector 
$\vec E(P)$, whose direction is given by the electric force 
$\vec F_{q_0}(P)$ on
a small, positive (so called ``test") charge $q_0$ placed at $P$, and whose
magnitude is given by $F_{q_0}(P)/q_0$:
%
\be
\vec E(P) = \vec F_{q_0}(P)/q_0
\ee
%
Note that by dividing by $q_0$, the electric field does not depend
on properties of the test charge.
It depends only on the charges producing the field.

\i For a single point charge $q$:
%
\be
E(P) = \frac{k q}{r^2}
\ee
%
where point $P$ is a distance $r$ from $q$.
The direction of the electric field is radially away from the 
charge if $q>0$ and towards the charge if $q<0$.

\i The electric field produced by a set of discrete 
point charges $q_1$, $q_2$, $\cdots$, $q_N$ is given by the
superposition principle:
%
\be
\vec E(P) = \vec E_1(P) + \vec E_2(P) + \cdots + \vec E_N(P)
\ee
%
where $\vec E_1(P)$ is the electric field due to $q_1$, etc.

\ei

\subsection{Electric field examples}
\bi

\i Electric dipole: Two charges, $+q$ and $-q$ separated by distance
$a$ along the $x$-axis (located at $x=-a/2$, $x=a/2$, respectively).
Then the electric field at point $P=(0,y)$ on the $y$-axis is given by:
\be
E_y(P)=0\,,\qquad
E_x(P) = \frac{k q a}{\left[y^2 + (a/2)^2\right]^{3/2}}
\ee

When $y\gg a$, the electric field has magnitude $E(P)\simeq k q a/y^3$, 
which falls off as the {\it cube} of the distance to the dipole.

\i The electric field produced by an infinite 2-dimensional sheet of 
charge with surface charge density $\sigma$ has magnitude
%
\be
E(P) = \frac{|\sigma|}{2\epsilon_0}
\ee
%
The field points perpendicular to the sheet (away from the sheet if 
$\sigma>0$; toward the sheet if $\sigma<0$).
The magnitude of the electic field is independent of the distance from the sheet.

\i In practice, one doesn't have 2-dimensional sheets that are infinite in extent.
But the above results still hold for a finite sheet provided the point $P$ is 
sufficiently close to the sheet.

\i The electric field between two infinite, parallel sheets of charge with 
surface charge densities $+\sigma$ and $-\sigma$ has magnitude
%
\be
E(P) = \frac{\sigma}{\epsilon_0}
\ee
%
The electric field points from the positively-charged sheet to the 
negatively-charged sheet.

\i This last example will be important when we discuss 
{\it capacitors} later on.

\ei

%%%%%%%%%%%%%%%%%%%%%%%%%%%%%%%%%%%%5
\subsection{Electric field lines}

\bi

\i Pictorial representation of the electric field $\vec E$.  
Direction of field lines point in the direction of $\vec E$; 
density of field lines (number per area perpendicular to $\vec E$) 
is proportional to the magnitude of $\vec E$.

\i Rules for drawing electric field lines: 

(i) lines begin on $+$ charges, end on $-$ charges (or may 
extend to infinity);

(ii) the number of lines leaving or terminating on a charge
should be proportional to the magnitude of the charge;

(iii) electric field lines cannot cross.

\i For a point charge, the density of lines at a distance
$r$ from the charge is given by $N/4\pi r^2$, which is 
proportional to the magnitude of the field $E(P)=k q/r^2$
as it should be.
 
\ei

%%%%%%%%%%%%%%%%%%%%%%%%%%%%%%%%
\subsection{Conductors in electrostatic fields}
\bi

\i Properties:

1) $\vec E=\vec 0$ everywhere inside a conductor, for both
solid conductors and conductors with {\em empty} cavities.

2) If an isolated conductor carries a charge, then that excess
charge resides on the outer surface of the conductor.

3) The electric field just outside the surface of a charged 
conductor is perpendicular to the surface with magnitude
$E= \sigma/\epsilon_0$, where $\sigma$ is the surface charge
density.

4) For irregularly shaped conductors, the surface charge 
density $\sigma$ is greatest where the radius of curvature of
the surface is smallest (i.e., where the surface is sharpest).

\i Property (1) arises since charges in the conductor 
rearrange themselves in response to an applied electric 
field such that the {\em induced} field associated with 
these charges exactly cancels the applied field in the 
conductor, $\vec E+\vec E_{\rm ind}=\vec 0$.

\i Property (1) implies that any external electric field
is shielded from a cavity inside a conductor.
That's why you're safe inside a car during a lightning
storm.

\i Fields within a cavity, produced e.g., by a charge,
do make it outside of the conductor.  So the outside of
a conductor is not shielded from the fields inside a 
cavity of a conductor.

\ei

%%%%%%%%%%%%%%%%%%%%%%%%%%
\end{document}


\section{Electrostatic potential}

\subsection{Work and potential energy}

\bi

\i Recall: The work done by a force in moving a particle
from point $A$ to point $B$ along some path is given
by 
%
\be
W_{A\rightarrow B} = \int_A^B \vec F\cdot d\vec s
\ee
%where $d\vec s$ is the infinitesimal displacement vector
%along the path.

\i For a general force, the value of the integral depends
on the path connecting $A$ and $B$.
If the force is {\em conservative}, then the integral
is path independent and depends only on the endpoints.

\i Examples of conservative forces are gravity, the 
spring force, and the electrostatic force.
An example of a non-conservative force is friction.

\i For conservative forces, one can a associate a 
{\em potential energy} $U$ with the work done by 
the force in moving the particle from $A$ to $B$:
%
\be
\Delta U\equiv U_B- U_A = -W_{A\rightarrow B} = 
-\int_A^B \vec F\cdot d\vec s
\ee

\i Note the minus sign in the above expression.
The potential energy of the system decreases if the 
force does positive work.
For example, as a mass $m$ falls in a gravitational
field $\vec g$, it loses potential energy, while gaining
kinetic energy.

\ei

%%%%%%%%%%%%%%%%%%%%%%%%%%%%%%%%
\subsection{Electrostatic potential}

\bi

\i The above results apply to the electrostatic force 
acting on a charge $q$.  
We define
%
\be
\Delta U = -\int_A^B \vec F_e\cdot d\vec s
= -q\int_A^B \vec E\cdot d\vec s
\ee

\i The electrostatic potential $V$ is then defined as 
the electrostatic potential energy per unit charge:
\be
\Delta V \equiv \frac{\Delta U}{q}
= -\int_A^B \vec E\cdot d\vec s
\ee
%
which has units of Joules per Coulomb, or volts
($1~{\rm volt} \equiv 1~{\rm J}/{\rm C}$).


\i Note that changing the electrostatic potential
$V$ by an additive constant ($V'= V+{\rm const}$) 
does not change the potential difference, since 
$\Delta V = \Delta V'$.
So the zero (or reference point) of the 
electrostatic potential can be set in whatever way
is most convenient for the problem at hand.
%(We will see that for point charges, the reference
%point is taken to be infinity.)

\i If $\vec E$ is everywhere orthogonal to $d\vec s$,
then $\Delta V=0$.
This means that the electric field $\vec E$ is 
{\em perpendicular} to an
{\em equipotential} surface.

\ei
%%%%%%%%%%%%%%%%
\subsection{Obtaining $\vec E$ from $V$}
\bi

\i The electrostatic potential $V$ is given by integrating
the electric field $\vec E$.
Conversely, by considering two nearby points $A$ and $B$, 
one can 
show that the electric field $\vec E$ is given by 
differentiating $V$:
%
\be
E_x = -\frac{\partial V}{\partial x}\,,\quad
E_y = -\frac{\partial V}{\partial y}\,,\quad
E_z = -\frac{\partial V}{\partial z}
\quad\Leftrightarrow\quad
\vec E=-\grad V
\ee
%

\i This relationship is very useful since it is often much
easier to first calculate the electrostatic potential $V$
for a distribution of charges, and then differentiate $V$
to obtain the electric field $\vec E$.

\i Recall that it is typically hard to calculate $\vec E$
for a distribution of charges since one has to sum up
{\em vectors} for each charge in the distribution.
Summing up scalars (for the potential $V$) is much easier.

\i Note that the units of $\vec E$ can also be expressed
in terms of volt per meter:
$1~{\rm N}/{\rm C} = 1~{\rm volt}/{\rm m}$.

\ei

\end{document}

